\section{Cosmology}
\label{sec:cosmo}

Hypothesis H3 says that the superfluid explains the dark matter, or the Hubble tension, or both. The
earlier paper examined it with the extra coordinates static. It found that the dark radiation of the
model does not relieve the Hubble tension and that a superfluid of the fermions cannot be the dark
matter of galaxies, although heavy bound pairs could be cold dark matter of an ordinary kind
\cite{\SIG{Flat1644}{Flat1671}}. This section examines what changes when $x5,x6,x7$ deflate.
The deflation enters cosmology in several ways, and each requires an assumption beyond the
hypothesis of deflation. The assumptions are stated first and are kept apart throughout, so that the
reader can see which conclusion rests on which.

\subsection{Assumptions, and the equation of the expansion}

\emph{A1: ordinary space expands.} The coordinates $x1,x2,x3$ carry the scale factor
$a_{\mathrm{sp}}(x4)$ of standard cosmology, with the Hubble rate $H=d\ln a_{\mathrm{sp}}/dx4$; the
hidden coordinate $x0$ is static; $x5,x6,x7$ deflate at constant rates. The metric of
Sec.~\ref{sec:geometry} is the case $a_{\mathrm{sp}}=1$.

\emph{A2: gravity is the Einstein theory in eight dimensions}, with a constant $G_8$. The constant
of gravitation that an observer in ordinary space measures is then $G_8$ divided by the proper volume
of the extra coordinates,
\begin{equation}
\GN(x4)\ \propto\ \frac{1}{\Vd}\ \propto\ e^{\gsum\,x4},\qquad\frac{1}{\GN}\frac{d\GN}{dx4}=\gsum .
\label{eq:GN}
\end{equation}
Newton's constant \emph{grows} as the extra volume shrinks: gravity was weaker in the past. That a
changing volume of extra dimensions makes the constants of four dimensions change, and that
nucleosynthesis bounds the change, was pointed out long ago \cite{KolbPerryWalker1986}; for a review
of varying constants see Ref.~\cite{Uzan2011}. By Sec.~\ref{sec:egb}, A2 is an assumption that the
deflating spacetime itself puts in doubt: without matter the metric solves the equations of
Einstein-Gauss-Bonnet gravity, not Einstein's.

For a diagonal metric whose seven directions other than the time have the rates $H,H,H$ (ordinary
space), $0$ (the coordinate $x0$) and $-\gff,-\gss,-\gvv$, the component of Einstein's equations
along the time is the sum of the products of all pairs of rates, set equal to $8\pi G_8$ times the
energy density in eight dimensions. Dividing the density by the extra volume turns $G_8$ into $\GN$:
\begin{equation}
3H^{2}-3\,\gsum\,H+\gprod=8\pi\,\GN(x4)\,\rho,
\label{eq:friedmann}
\end{equation}
with the density $\rho$ of radiation, matter and the cosmological term in ordinary space. Equation
(\ref{eq:friedmann}) is the only one of Einstein's equations that is used; the others would
determine the stresses of the unknown source that keeps the three rates constant.

The calculations take three equal rates, so that $\gprod=\gsum^{2}/3$, and the densities of the
standard model of cosmology \cite{Planck2018VI} as the earlier paper used them: baryons
$\omega_b=0.02237$, cold dark matter $\omega_c=0.1200$, and $H_0=67.36~\kmsMpc$ without deflation.
The cosmological term is chosen so that Eq.~(\ref{eq:friedmann}) holds today. Without deflation the
program reproduces the two numbers of the earlier calculation on which this section rests: the angle
$\theta_*$ under which the sound horizon at last scattering is seen, $100\,\theta_*=1.03965$, and
that sound horizon, $r_s=144.449$~Mpc.

\subsection{The Hubble constant inferred from the sound horizon}

The microwave background fixes the angle $\theta_*$ to high precision. A model that changes the history of the expansion must keep that angle,
and the Hubble constant that it then requires is the \emph{inferred} Hubble constant. The tension is
that the standard model infers $67.4~\kmsMpc$ \cite{Planck2018VI} while the local distance network
measures $73.50\pm0.81$ \cite{H0DN2026}. Table~\ref{tab:hubble} gives the inferred value on the
deflating spacetime under A1 and A2, at fixed $\theta_*$ and fixed densities of baryons and dark
matter.

\begin{table*}
\caption{\label{tab:hubble}The Hubble constant inferred from the angular scale of the sound horizon
if Newton's constant grows with the deflation (assumptions A1 and A2). $\GN$ ratio: Newton's constant
in the early universe divided by its value today, $e^{-\gsum\,t_0}$ with the age $t_0$ of the model.
The rows at $\gsum/H_0=0.00324$ and $0.0644$ are the bounds of lunar laser ranging and of
nucleosynthesis (Table~\ref{tab:bounds}).}
\begin{ruledtabular}
\begin{tabular}{ccccccc}
$\gsum/H_0$ & $\gsum$ (per year) & $H_0$ ($\kmsMpc$) & change & $\GN$ ratio & $r_s$ (Mpc) & $t_0$ (Gyr) \\
\hline
{{ROWS:tbl_hubble}}
\end{tabular}
\end{ruledtabular}
\end{table*}

\emph{The inferred Hubble constant falls}, with the slope $d\ln H_0/d(\gsum/H_0)=-0.911$ at small
rates. The reason is this. By Eq.~(\ref{eq:GN}) gravity was weaker in the early universe; the
universe expanded more slowly at a given density; sound had more time to travel, and the sound
horizon is larger. To keep the angle, the distance to last scattering must be larger, which requires
a smaller Hubble constant. The deflation moves the inferred value \emph{away} from the local
measurement [Fig.~\ref{fig:hubble}(a)]. The term $\gprod$ of Eq.~(\ref{eq:friedmann}) plays no part
in this: it is of the second order in the rate, and it acts as a constant, which the cosmological
term that is adjusted to today's expansion absorbs almost exactly.

\begin{figure*}
\includegraphics[width=\textwidth]{fig_hubble}
\caption{\label{fig:hubble}The Hubble constant and the growth of structure on a deflating spacetime.
(a) The Hubble constant inferred from the angular scale of the sound horizon against the rate of
deflation, if gravity feels the deflation (A1, A2); the bounds of Table~\ref{tab:bounds} and the
local measurement are marked. (b) The change of the inferred Hubble constant if Newton's constant is
held fixed and bound pairs lose rest mass. (c) Amplitude and rate of the growth of structure relative
to the model without deflation.}
\end{figure*}

\subsection{How fast can the coordinates deflate today?}

\begin{table*}
\caption{\label{tab:bounds}Upper bounds on the present total rate of deflation, each under the
assumption named. Uncertainties of the measured values are at 68 per cent unless stated; the bounds
are at two standard deviations. The bound in units of $H_0$ uses $H_0=67.36~\kmsMpc$. A2: gravity is
the Einstein theory in eight dimensions. HV: the vacuum value of the Vorton-Scalar sets the
electroweak scale. H2: the superfluid state is the Higgs field; the bound is given for $\nEW=1$ and
scales as $1/\nEW$ (Sec.~\ref{sec:higgs}).}
\begin{ruledtabular}
\begin{tabular}{p{0.17\textwidth}p{0.27\textwidth}p{0.15\textwidth}cc}
observation & measured & assumption & bound on $\gsum$ (per year) & in units of $H_0$ \\
\hline
lunar laser ranging \cite{HofmannMueller2018} & $(d\GN/dt)/\GN=(7.1\pm7.6)\times10^{-14}$ per year & A2 & {{ROWS:tbl_bounds|rows=0|cols=2,3}}
nucleosynthesis, Newton's constant \cite{Alvey2020} & $\GN$ then, divided by $\GN$ today: $0.99^{+0.06}_{-0.05}$ (two standard deviations) & A2; constant rate & {{ROWS:tbl_bounds|rows=1|cols=2,3}}
nucleosynthesis, Fermi constant \cite{ScherrerSpergel1993} & $G_{\mathrm{Fermi}}$ at $t=1$~s within 20 (10) per cent of today's value, at fixed fermion masses & HV; constant rate & {{ROWS:tbl_bounds|rows=2|cols=2,3}}
atomic clocks, mass ratio \cite{Lange2021} & $(d/dt)\ln(m_p/m_e)=(-8\pm36)\times10^{-18}$ per year & H2, $\nEW=1$ & {{ROWS:tbl_bounds|rows=3|cols=2,3}}
atomic clocks, fine-structure constant \cite{Lange2021} & $(d\alpha_{\mathrm{em}}/dt)/\alpha_{\mathrm{em}}=(1.0\pm1.1)\times10^{-18}$ per year & an electromagnetic field along $x5,x6,x7$ (not part of the model) & {{ROWS:tbl_bounds|rows=4|cols=2,3}}
\end{tabular}
\end{ruledtabular}
\end{table*}

Table~\ref{tab:bounds} and Fig.~\ref{fig:bounds} collect the bounds.

\emph{Under A2, lunar laser ranging} measures $(d\GN/dt)/\GN$ directly, and Eq.~(\ref{eq:GN}) turns
its limit into $\gsum<2.23\times10^{-13}$ per year, that is, $3.24\times10^{-3}\,H_0$. At this rate
the deflation clock has advanced by $\gsum\,t_0=3.1\times10^{-3}$ over the age $t_0$ of the universe,
and the inferred Hubble constant is lower by $0.20~\kmsMpc$.

\emph{Under A2, nucleosynthesis} bounds the ratio of Newton's constant then to its value today; for
a rate that was constant since, $\gsum<0.064\,H_0$. This is weaker than the first bound but does not
assume that the rate today equals the rate in the past.

\emph{Under HV, nucleosynthesis} bounds the change of the Fermi constant at fixed fermion masses to
10 to 20 per cent, which gives $\gsum<0.111$ to $0.234\,H_0$. HV is the assumption that the
Vorton-Scalar itself plays the role of the Higgs field, so that
$G_{\mathrm{Fermi}}=1/(\sqrt2\,\vVS^{2})$ grows as $e^{\td}$ while every mass stays constant. It is
\emph{not} the hypothesis H2, and the earlier paper already noted that the Vorton-Scalar, a singlet
of all non-Abelian symmetries, cannot be the Higgs field of the Standard Model; the bound is given
only to show what the assumption would imply.

\emph{Under the hypothesis H2, atomic clocks} give by far the strongest bound, and one that needs no
assumption about gravity. Section~\ref{sec:higgs} derives it: $\gsum<1.16\times10^{-6}\,H_0/\nEW$
with an exponent $\nEW$ of order one or larger.

\emph{The fine-structure constant} is bounded by the same clocks to parts in $10^{18}$ per year.
This would bound the rate, to $4.6\times10^{-8}\,H_0$, only if an electromagnetic field propagated
along the deflating coordinates. The model contains no electromagnetic field; the row shows how
strongly such an extension would be constrained.

\emph{In the model as it stands}, without H2, HV or an electromagnetic extension and without A2, no
mass and no coupling of ordinary matter depends on the deflation, and none of these observations
bounds the rate.

The consequence for the histories of Secs.~\ref{sec:gap} to \ref{sec:cjt} is severe. Under A2 the
clock has advanced by at most 0.0031 during the whole of cosmic history, if the rate was constant;
under H2 by at most $1.1\times10^{-6}/\nEW$. A history in which pairs become bound, matter loses its
fermion mass or the gap grows by orders of magnitude within the age of the universe requires that
gravity does not feel the deflation, that the superfluid is not the Higgs field, or that the rate
was much larger in the past than it is today.

\begin{figure*}
\includegraphics[width=0.8\textwidth]{fig_bounds}
\caption{\label{fig:bounds}Upper bounds on the present total rate of deflation in units of the
Hubble constant, each under the assumption of Table~\ref{tab:bounds}. The allowed region is to the
left of each bound.}
\end{figure*}

\subsection{Bound pairs lose rest mass}

If the dark matter consists of bound pairs of the fermions, its particles have the rest mass $2m$
minus the binding energy. The binding energy grows with the coupling, so that the pairs become
lighter as the deflation proceeds, and the density of dark matter falls slightly faster than the
inverse volume. A dark matter that was heavier in the past raises the inferred Hubble constant. The
program computes the effect for four examples with $\aVS$ between 0.4 and 1 today, where the binding
energy is at most 8.2 per cent of the rest mass [Fig.~\ref{fig:hubble}(b)]. If gravity feels the
deflation, the inferred Hubble constant changes by between $-4.11$ and $-0.19~\kmsMpc$ at the rates
that nucleosynthesis and lunar laser ranging allow: the weaker gravity of the past outweighs the
loss of mass. If Newton's constant is held fixed and the deflation acts on the dark sector alone,
rates between 0.03 and $1\,H_0$ raise the inferred Hubble constant by at most $+0.48~\kmsMpc$, for a
loss of up to 7.7 per cent of the rest mass. This is the only effect of the deflation found here
that raises the inferred Hubble constant, and in these examples it supplies less than a tenth of the
$6.1~\kmsMpc$ that separate the model from the local measurement.

\subsection{The growth of structure}

Under A1 and A2 the density contrast $\delta_m$ of matter obeys, in the Newtonian limit, the usual
equation with the Newton's constant of the moment,
\begin{equation}
\frac{d^{2}\delta_m}{dx4^{2}}+2H\,\frac{d\delta_m}{dx4}=4\pi\,\GN(x4)\,\rho_m\,\delta_m ,
\end{equation}
with the density $\rho_m$ of matter. Two effects oppose each other [Fig.~\ref{fig:hubble}(c)]. At a
\emph{fixed} Hubble constant, weaker gravity in the past lowers the amplitude $\sigma_8$ of clustering
today: $d\ln\sigma_8/d(\gsum/H_0)=-0.317$. But the Hubble constant that keeps the angle of the sound
horizon is lower, the cosmological term takes over later, and matter clusters for longer; with that
Hubble constant the net change is $+0.102$ per unit of $\gsum/H_0$. The amplitude therefore barely
moves, by $+0.03$ per cent at the bound of lunar laser ranging and $+0.67$ per cent at the bound of
nucleosynthesis, while the growth rate today, $f_g=d\ln\delta_m/d\ln a_{\mathrm{sp}}$, rises by the
factors 1.0027 and 1.0589. A deflating universe would show itself in the product $f_g\sigma_8$,
which redshift-space distortions measure, rather than in $\sigma_8$.

\subsection{Dark matter}
\label{sec:dm}

\emph{What does not change.} The estimates of the earlier paper that decide whether the particles of
the model can be degenerate or condensed in observed halos contain the masses, the densities and the
velocities only, and the masses do not change along the deflation. Computed again, every number
agrees: fermions with eight internal states are degenerate in the core of a dwarf galaxy only if
they are lighter than 0.276~keV, and bound pairs are condensed there only if they are lighter than
262~eV. The verdict of the earlier paper on these two possibilities stands.

\emph{Self-interaction, and the epoch of the resonance.} The cross-section per mass of bound pairs
must respect the bound of about $1~\mathrm{cm^{2}/g}$ from colliding clusters
\cite{Markevitch2004,Randall2008}, which the earlier paper turned into a smallest allowed fermion
mass: 918~MeV for $\aVS=0.1$. For a given particle the masses are constant and so is this
cross-section. What the deflation changes is whether the pair \emph{exists}: it is bound only after
$\DVS$ has grown to $D_c$. And at that moment the scattering length of two unbound fermions
diverges (Sec.~\ref{sec:pair}), so that their cross-section rises to the limit that unitarity
allows, which is larger the slower the particles. For the example particle, $m=918$~MeV and
$m_s=54.7$~MeV, at the relative velocities 10, 156 and 1000~km/s [Fig.~\ref{fig:binding}(c)]: one
unit of the clock before the binding the cross-section per mass is $0.84~\mathrm{cm^{2}/g}$ at all
three velocities; at the binding it peaks at $1.3\times10^{7}$, $5.2\times10^{4}$ and
$1.3\times10^{3}~\mathrm{cm^{2}/g}$; and it exceeds $100~\mathrm{cm^{2}/g}$ during an interval of the
clock of 0.27, 0.27 and 0.26 around the binding. \emph{A deflating universe whose dark matter
consists of these fermions passes through an epoch of strong, velocity-dependent self-interaction},
strongest in the smallest halos. Whether the epoch lies in the observable past depends on the rate.
Under A2 and a constant rate it does so only if $\DVS$ today exceeds $D_c$ by less than a third of a
per cent, a coincidence that nothing in the model explains.

\emph{How fast the quantities drift.} With $n=d\ln(\text{quantity})/d\td$, the relative change of a
quantity over the age of the universe is $n\,\gsum\,t_0$. The couplings themselves drift with $n=1$
or $\tfrac12$ (Table~\ref{tab:exponents}). Quantities near a threshold drift much faster: the
scattering length just below the binding, at $\DVS=1.5$, with $n=9.1$; the binding energy just above
it, at $\DVS=1.8$, with $n=30$; the gap of a weakly paired state with $n=1/\lambda$. These are the
quantities in which a slow deflation would first be seen.

%%<1671>
\subsection{The quanta that the deflation creates}

The temperature (\ref{eq:Tdefl}) of the quanta that the deflation creates in modes that depend on a
deflating coordinate is, for three equal rates at the bound of lunar laser ranging,
$2.9\times10^{-33}$~K. The number of created particles of mass $m$ per mode is $e^{-2\pi m/\rt}$, with an
exponent above $10^{37}$ for every mass above 1~keV at any rate allowed today. \emph{No particle is
created today.} Creation by deflation matters only for an epoch in which the rate was comparable
with the masses, which no bound of this section addresses.
%%</1671>

\subsection{Verdict on hypothesis H3}

\emph{The Hubble tension.} The mechanism of the earlier paper, the phase mode of the Vorton-Scalar
as dark radiation, is untouched at the allowed rates, since the couplings drift by parts in a
thousand; its verdict stands and is not improved. The deflation itself contributes two effects. The
first, through Newton's constant, has the wrong sign and is bounded to $0.2~\kmsMpc$. The second,
the loss of mass of bound pairs, has the right sign and is bounded, in the examples computed, to
$0.5~\kmsMpc$. \emph{The hypothesis that the model resolves the Hubble tension is not supported by
the deflation, and under A2 it is mildly disfavored.}

\emph{Dark matter.} The conclusions of the earlier paper about the mass ranges hold unchanged. The
deflation adds one qualitative prediction, the epoch of the resonance, and one quantitative one, the
shift of the growth rate. Neither is in conflict with observation at the allowed rates, and neither
is detectable at the rate that lunar laser ranging allows. \emph{The hypothesis that the superfluid
is the dark matter is neither strengthened nor weakened by the deflation.} It acquires a history,
whose observable part is empty if gravity is eight-dimensional and the rate constant.

\emph{What was not computed.} The earlier paper computed the temperature and polarization spectra
of the microwave background with its own Einstein-Boltzmann program and forecast the constraints
from them. That calculation has not been repeated with a Newton's constant that depends on time,
for two reasons. First, the equations of the perturbations contain the stresses of the source that
keeps the rates of deflation constant, which the hypothesis does not specify; a spectrum can be
computed only under a further assumption about that source or about the theory of gravity. Second,
it requires a version of the program in which the gravitational coupling of the background and of
the perturbations depends on time while the rates of Thomson scattering and of recombination do not,
validated again against the constant case. What is established here is the effect on the
\emph{background}. At the bound of lunar laser ranging Newton's constant at last scattering was
0.9969 of today's value; the effect on the spectra beyond the shift of the angle of the sound
horizon is of that relative size and has not been computed.

